% SEGMENT OLP-0034-S01
% Part:sets-functions-relations
% Chapter: sets
% Section: reduction

\documentclass[../../../include/open-logic-section]{subfiles}

\begin{document}

\olfileid{sfr}{siz}{red}

\olsection{குறைத்தல்}

\begin{editorial}
இந்தப் பிரிவு, \olref[nen]{sec} இல் உள்ள முடிவுகளுக்கு இணையாக,
குறைத்தல் மூலம் எண்ணிடத்தகாத தன்மையை நிறுவுகிறது.
\olref[nen-alt]{sec} இல் உள்ள முடிவுகளுக்கு இணையான,
சற்று மேலும் சுருக்கமான மாற்று வடிவம்
\olref[red-alt]{sec} இல் கொடுக்கப்பட்டுள்ளது.
\end{editorial}

$\Pow{\PosInt}$ என்பது !!{nonenumerable} என ஒரு
மூலைவிட்டமாக்க வாதத்தின் மூலம் காட்டினோம்.
எல்லா முடிவுறாத $0$, $1$ தொடர்களின் கணமான $\Bin^\omega$
என்பது !!{nonenumerable} என்பதற்கான நிறுவல் ஏற்கெனவே நம்மிடம் இருந்தது.
$\Pow{\PosInt}$ என்பது !!{nonenumerable} என நிறுவுவதற்கான
மற்றொரு வழி இதோ:
\emph{$\Pow{\PosInt}$ என்பது !!{enumerable} எனில்,
$\Bin^\omega$ உம் !!{enumerable}} எனக் காட்டுக.
$\Bin^\omega$ என்பது !!{enumerable} அல்ல என்பது நமக்குத் தெரிவதால்,
$\Pow{\PosInt}$ உம் அவ்வாறு இருக்க முடியாது.
இது ஒரு சிக்கலை மற்றொரு சிக்கலாகக் \emph{குறைத்தல்} எனப்படுகிறது.
இங்கு $\Bin^\omega$ ஐப் பட்டியலாக்கும் சிக்கலை,
$\Pow{\PosInt}$ ஐப் பட்டியலாக்கும் சிக்கலாகக் குறைக்கிறோம்.
பிந்தையதற்கான ஒரு தீர்வு, அதாவது $\Pow{\PosInt}$ இன் பட்டியலாக்கம்,
முந்தையதற்கான ஒரு தீர்வை, அதாவது $\Bin^\omega$ இன் பட்டியலாக்கத்தை,
வழங்கும்.

ஒரு கணம் $B$ ஐப் பட்டியலாக்கும் சிக்கலை,
ஒரு கணம் $A$ ஐப் பட்டியலாக்கும் சிக்கலாக எவ்வாறு குறைப்பது?
$A$ இன் ஒரு பட்டியலாக்கத்தை $B$ இன் பட்டியலாக்கமாக மாற்றும்
ஒரு வழியை வழங்குகிறோம்.
இதற்கான எளிய வழி, !!a{surjective} சார்பு
$f\colon A \to B$ ஐ வரையறுப்பதாகும்.
$x_1$, $x_2$, \dots{} என்பது $A$ ஐப் பட்டியலாக்கினால்,
$f(x_1)$, $f(x_2)$, \dots{} என்பது $B$ ஐப் பட்டியலாக்கும்.
நமது நிலையில், ஒரு மேற்கோர்த்தல் சார்பு
$f\colon \Pow{\PosInt} \to \Bin^\omega$ ஐத் தேடுகிறோம்.

% SEGMENT OLP-0034-S02
\begin{prob}
!!{injective} சார்பு $g\colon B \to A$ ஒன்று இருந்தும்,
$B$ என்பது !!{nonenumerable} ஆகவும் இருந்தால்,
$A$ உம் அவ்வாறே எனக் காட்டுக.
$A$ இன் ஒரு பட்டியலாக்கத்தை $B$ இன் பட்டியலாக்கமாக மாற்றுவதற்கு
$g$ ஐ எவ்வாறு பயன்படுத்தலாம் எனக் காட்டி இதைச் செய்க.
\end{prob}

% SEGMENT OLP-0034-S03
\begin{proof}[\olref[nen]{thm:nonenum-pownat} இன் குறைத்தல் நிறுவல்]
$\Pow{\PosInt}$ என்பது !!{enumerable} எனவும், எனவே அதன்
$Z_{1}$, $Z_{2}$, $Z_{3}$, \dots என்ற ஒரு பட்டியலாக்கம்
இருப்பதாகவும் கொள்வோம்.

% SOURCE CORRECTION OLSIZ-004; see translation/size-notes.tex.
$f \colon \Pow{\PosInt} \to \Bin^\omega$ என்ற சார்பை,
$f(Z)$ என்பது பின்வரும் $s$ தொடராக இருக்குமாறு வரையறுப்போம்:
$n \in Z$ ஆக இருக்கும்போது, அப்போதும் அப்போது மட்டுமே,
$s(n) = 1$; மற்ற நிலைகளில் $s(n) = 0$.
இது தெளிவாக ஒரு சார்பை வரையறுக்கிறது.
ஏனெனில் $Z \subseteq \PosInt$ எனும்போது,
எந்த $n \in \PosInt$ உம் $Z$ இன் !!a{element} ஆக இருக்கிறது
அல்லது அதன் உறுப்பு அல்ல.
எடுத்துக்காட்டாக, நேர்ம இரட்டை எண்களின் கணமான
$2\PosInt = \{2,
4, 6, \dots\}$ என்பது $010101\dots$ என்ற தொடருக்கும்,
வெற்றுக் கணம் $0000\dots$ என்ற தொடருக்கும்,
$\PosInt$ என்ற கணமே $1111\dots$ என்ற தொடருக்கும் அனுப்பப்படுகிறது.

இது !!{surjective} பண்பையும் கொண்டுள்ளது:
$0$, $1$ ஆகியவற்றின் ஒவ்வொரு தொடரும், ஏதாவது ஒரு நேர்ம முழுக்களின்
கணத்திற்கு ஒத்திருக்கிறது.
அந்தக் கணத்தின் உறுப்புகள், தொடரில் $1$ உள்ள இடங்களுக்கு
ஒத்த முழுக்களாகும்.
மேலும் துல்லியமாக, $s \in \Bin^\omega$ எனக் கொள்வோம்.
$Z \subseteq \PosInt$ ஐப் பின்வருமாறு வரையறுப்போம்:
\[
Z = \Setabs{n \in \PosInt}{s(n) = 1}
\]
அப்போது $f$ இன் வரையறையைப் பார்த்து,
$f(Z) = s$ என்பதைச் சரிபார்க்கலாம்.

இப்போது பின்வரும் பட்டியலைக் கருதுக:
\[
f(Z_1), f(Z_2), f(Z_3), \dots
\]
$f$ என்பது !!{surjective} பண்புடையதால்,
$\Bin^\omega$ இன் ஒவ்வோர் உறுப்பும் ஏதாவது ஓர் உள்ளீட்டிற்கு
$f$ இன் மதிப்பாகவும், எனவே பட்டியலில் இடம்பெறவும் வேண்டும்.
ஆகவே இந்தப் பட்டியல் $\Bin^\omega$ முழுவதையும் பட்டியலாக்க வேண்டும்.

எனவே $\Pow{\PosInt}$ என்பது !!{enumerable} எனில்,
$\Bin^\omega$ உம் !!{enumerable} ஆகும்.
ஆனால் $\Bin^\omega$ என்பது !!{nonenumerable}
(\olref[nen]{thm:nonenum-bin-omega}).
எனவே $\Pow{\PosInt}$ என்பது !!{nonenumerable}.
\end{proof}

% SEGMENT OLP-0034-S04
\begin{explain}
குறைத்தல் எந்தத் திசையில் செல்கிறது என்பதில் குழப்பமடைவது எளிது.
எடுத்துக்காட்டாக, !!a{surjective} சார்பு
$g \colon \Bin^\omega \to B$ இருப்பது,
$B$ என்பது !!{nonenumerable} என \emph{நிறுவாது}.
($g(s) = s(1)$ என வரையறுக்கப்படும்
$g \colon \Bin^\omega \to \Bin$ ஐக் கருதுக.
இந்தச் சார்பு $0$, $1$ தொடர்களை அவற்றின் முதல் !!{element} க்கு
அனுப்புகிறது.
சில தொடர்கள் $0$ இலும் சில தொடர்கள் $1$ இலும் தொடங்குவதால்,
இது !!{surjective} பண்புடையது.
ஆனால் $\Bin$ முடிவுறு கணமாகும்.)
$f$ என்ற சார்பு !!{surjective} பண்புடையதாக இருக்க வேண்டும்
என்பதையும் கவனிக்கவும்; இல்லையெனில் வாதம் செயல்படாது.
அவ்வாறு இல்லையெனில், $f(x_1)$, $f(x_2)$, \dots{} என்பது
$B$ இன் எல்லா !!{element}s[acc]யும் கொண்டிருப்பது உறுதியாகாது.
எடுத்துக்காட்டாக,
% SOURCE CORRECTION OLSIZ-005; see translation/size-notes.tex.
\[
h(n) = \underbrace{000\dots0}_{\text{$n$ $0$கள்}}111\dots
\]
என்பது $h\colon \PosInt \to
\Bin^\omega$ என்ற சார்பை வரையறுக்கிறது;
ஆனால் $\PosInt$ என்பது !!{enumerable}.
\end{explain}

% SEGMENT OLP-0034-S05
\begin{prob}
நேர்ம முழுக்களின் ஜோடிகளைக் கொண்ட எல்லாக்
\emph{கணங்களையும்} கொண்ட கணம் !!{nonenumerable} என்பதை
ஒரு குறைத்தல் வாதத்தின் மூலம் காட்டுக.
\end{prob}

% SEGMENT OLP-0034-S06
\begin{prob}\label{sfr:siz:red:prob:nat-nat}
$f\colon \Nat \to \Nat$ என அமையும் எல்லாச் சார்புகளையும் கொண்ட
கணம் $X$ என்பது !!{nonenumerable} என்பதை ஒரு குறைத்தல் வாதத்தின்
மூலம் காட்டுக.
(குறிப்பு: $X$ இலிருந்து $\Bin^\omega$ க்குச் செல்லும்
ஒரு மேற்கோர்த்தல் சார்பைக் கொடுக்கவும்.)
\end{prob}

% SEGMENT OLP-0034-S07
\begin{prob}
இயல் எண்களின் எல்லா முடிவுறாத தொடர்களையும் கொண்ட
$\Nat^\omega$ என்ற கணம் !!{nonenumerable} என்பதை
ஒரு குறைத்தல் வாதத்தின் மூலம் காட்டுக.
\end{prob}

% SEGMENT OLP-0034-S08
\begin{prob}
நேர்ம முழுக்களின் கணத்திலிருந்து $\{0\}$ என்ற கணத்திற்குச் செல்லும்
சார்புகளின் கணம் $P$ ஆகவும்,
நேர்ம முழுக்களின் கணத்திலிருந்து $\{0\}$ க்குச் செல்லும்
\emph{பகுதிச்} சார்புகளின் கணம் $Q$ ஆகவும் இருக்கட்டும்.
$P$ என்பது !!{enumerable} எனவும், $Q$ அவ்வாறு இல்லை எனவும் காட்டுக.
(குறிப்பு: $\Bin^\omega$ ஐப் பட்டியலாக்கும் சிக்கலை,
$Q$ ஐப் பட்டியலாக்கும் சிக்கலாகக் குறைக்கவும்.)
\end{prob}

% SEGMENT OLP-0034-S09
\begin{prob}
நேர்ம முழுக்களின் கணத்திலிருந்து \{0,1\} என்ற கணத்திற்குச் செல்லும்
எல்லா !!{surjective} சார்புகளையும் கொண்ட கணம் $S$ ஆக இருக்கட்டும்.
அதாவது, எல்லா !!{surjective} சார்புகள்
$f\colon \PosInt \to \Bin$ உம் $S$ இல் உள்ளன.
$S$ என்பது !!{nonenumerable} எனக் காட்டுக.
\end{prob}

% SEGMENT OLP-0034-S10
\begin{prob}
எல்லா மெய்யெண்களையும் கொண்ட கணம் $\Real$ என்பது
!!{nonenumerable} எனக் காட்டுக.
\end{prob}

\end{document}
